\begin{abstract}
Evaluating a language model on a full benchmark is expensive, so a growing literature selects a few items using the item-level results of existing (source) models and predicts the new (target) model's full accuracy from its answers to those items only. We ask when this source-model information actually beats the simplest practice, scoring a uniformly random subset. On 13 binary item-level score matrices (Open LLM Leaderboard v2 and HELM) we compare eleven selector--estimator pipelines under two protocols: a random split of models into sources and targets, and a capability-extrapolation split where the targets are the strongest 25\% of models. All methods see the same sources, targets and budgets ($n\le 200$ items), their hyperparameters never use target scores, and we report paired bootstrap intervals. On random splits every source-informed method has a lower macro-average MAE than the random subset mean at $n{=}50$ (4.6 accuracy points; 3.0--3.4 for the IRT, low-rank, supervised-selection and design-based pipelines, 3.7 for plain stratification, 4.3 for anchor points, which are no better than random from $n{=}100$), and rank correlation improves mostly at small $n$. Under extrapolation this advantage reverses for the methods that rely on a fitted model of the sources (MAE 6.5--8.1 points at $n{=}50$, worse than random) and their Spearman correlation drops from 0.64 to 0.46--0.56. What remains reliable is a design-based residual correction of the model's prediction (the difference estimator of survey sampling; ablations show that random items suffice and that it, not stratification, protects the model), and stratifying items by source difficulty when the estimate is the plain mean: they lower MAE by 0.7--1.6 points at $n{=}50$ in both regimes and are never significantly worse than random in any benchmark, split and budget, with intervals that also cover the choice of target models. The correction costs accuracy when targets resemble the sources, and its edge over plain stratification shrinks from about 1.5 points at $n{=}10$ to a third of a point at $n{=}50$. Source-model data thus help a lot when new models resemble the sources, only modestly otherwise, and no method recovers the ordering of models closer than a few accuracy points.
\end{abstract}
